% !TEX root = ../../Extensions of IQC Duality to Partial Integral Equations.tex
\subsection{Notation}
$L_2^n[a,b]$ denotes the set of $\R^n$-valued, Lebesgue square-integrable
equivalence classes of functions on the spatial domain $[a,b]\subset\R$. The
space $\R L_2^{\mathbf{n}}[a,b]$ denotes the Cartesian product $\R^{n_1} \times
L_2^{n_2}[a,b]$, where $\mathbf{n} = \bmat{n_1 \\ n_2}$, with inner product
\begin{equation*}
\ip{\bmat{x_1 \\ \mathbf{x}_2}}{\bmat{y_1 \\ \mathbf{y}_2}}_{\R L_2}
:= x_1^\top y_1 + \ip{\mathbf{x}_2}{\mathbf{y}_2}_{L_2}.
\end{equation*} When clear from context, we occasionally omit the subscript and
simply write $\ip{\cdot}{\cdot} := \ip{\cdot}{\cdot}_{\R L_2}$. Since the only
spatial domain considered here is $[a,b]$, we typically omit the domain and
write $L_2^n$ or $\R L_2^{\mathbf{n}}$, further omitting dimensions when clear
from context. We omit dimensions when clear from context.
For suitably differentiable spatial states, $\dot{\mbf x}:=\partial_t\mbf x$.
We denote the induced norm for
bounded linear operators on $\RL$ by
\[
\|\mcl P\|
:=\sup_{0\neq\mbf x\in\RL^{\mbf n}}
\frac{\|\mcl P\mbf x\|_{\RL}}{\|\mbf x\|_{\RL}}.
\]
Our notational convention is to write
spatial states and profiles in bold (e.g., $\mbf{x}\in\RL$).
For an operator $\mcl A$, the adjoint $\mcl
 A^*$ is defined by
\[
\ip{\mathbf{x}}{\mcl A \mathbf{y}}_{\RL} = \ip{\mcl A^* \mathbf{x}}{\mathbf{y}}_{\RL},
\quad \forall\, \mathbf{x},\mathbf{y} \in \RL.
\]
Define the signal space
\begin{align*}
L_{[a,b]}^{\mbf n}:=\{&\mbf u:[0,\infty)\to\RL^{\mbf n}[a,b] \mid\\&\quad\|\mbf u\|_{L_{[a,b]}}^2:=\int_0^{\infty}\|\mbf u(t)\|_{\RL}^2\,dt<\infty\}, 
\end{align*}
with inner product $\ip{\mbf u}{\mbf v}_{L_{[a,b]}}:=\int_0^{\infty}\ip{\mbf u(t)}{\mbf v(t)}_{\RL}\,dt$. 
For any $\mbf y:[0,\infty)\to\RL^{\mbf n}[a,b]$,
define the truncation operator $(p_T\mbf y)(t):=\mbf y(t)$ 
for $0\leq t\leq T$ and $(p_T\mbf y)(t):=0$ for $t>T$,
the corresponding extended signal space is therefore,
\begin{align*}
L_{\mrm{e},[a,b]}^{\mbf n}:=\{\mbf v \mid p_T\mbf v\in L_{[a,b]}^{\mbf n} \;\forall T>0\}.
\end{align*}
Moreover, for $\tau \in \R$ define the shift operator $(s_\tau u)(t):=u(t-\tau)$. An 
operator $G:L_{\mrm{e},[a,b]}^{\mbf n}\to
L_{\mrm{e},[a,b]}^{\mbf m}$ is called
\begin{itemize}
\item \emph{Causal}, if $p_TG=p_TGp_T$ for all $T>0$.
\item \emph{Anticausal}, if $(I-p_T)G(I-p_T)=(I-p_T)G$ for all $T>0$.
\item \emph{Memoryless}, if $G$ is simultaneously causal and anticausal.
\item \emph{Time-Invariant}, if $s_\tau G = G s_\tau$ for all $\tau\in\R$.
\item \emph{Bounded on $L_{[a,b]}$}, if $\|G\|\leq C$ for some $C\geq0$.
\item \emph{Bounded on $L_{\mrm{e},[a,b]}$}, if
$\|p_TG\|\leq C$ for some $C\geq0$ and all $T>0$.
\end{itemize}
For a bounded linear operator
$\mcl K:\RL^{\mbf n}\to\RL^{\mbf m}$, define the multiplication operator
$\mcl M_{\mcl K}:L_{\mrm{e},[a,b]}^{\mbf n}
\to L_{\mrm{e},[a,b]}^{\mbf m}$ by
\[
(\mcl M_{\mcl K}u)(t):=\mcl K u(t).
\]

\subsection{Partial Integral Equations (PIEs)}
Partial Integral operators, also known as PI operators, are integral operators
defined jointly on a vector space and Hilbert
space~\cite{peet2021pierepresentation,shivakumar2024GDPE}.
\begin{definition}
We say an operator $\mcl P:\RL^{\mbf n}[a,b]\to\RL^{\mbf m}[a,b]$
belongs to ${\Pi}_4^{\mbf{m}\times\mbf{n}}$, denoted by
$\PIop{P}{Q_1}{Q_2}{\{R_0,R_1,R_2\}}$, if there exist a matrix $P$ and matrix polynomials $Q_1, Q_2,
R_0, R_1$, and $R_2$ (of compatible dimension) such that,
\begin{gather*}
\Bigg[\Bigg(\PIop{P}{Q_1}{Q_2}{\{R\}}\Bigg)\! \!
\bmat{x \\ \mbf{x}}\Bigg](s):=\! \bmat{Px\!
+\!\int\limits_a^b\!Q_1(s)\mbf x(s)ds\\Q_2(s)x\!+\!\mcl R \mbf x(s)},\\ 
(\mcl R \mbf x)(s)\!\! =\!\! R_0(s)\mbf x(s)\! +\! \int\limits_a^s\!\!\!R_1(s,\theta)\mbf x(\theta)d\theta\! +\!\int\limits_s^b\!\!\!R_2(s,\theta)\mbf x(\theta) d\theta.
\end{gather*}
\end{definition}
The vector space of PI operators, given compatible dimensions, are closed under
composition, addition, adjoint and concatenation -- implying that ${\Pi}_4$ is a
composition algebra \cite{shivakumar2024GDPE}.

\begin{definition}(LPIs)
\label{def:lpi-operator}
For given PI operators $\{\mcl E_{i,j}, \mcl{F}_{i,j}, \mcl{G}_i\}$ and a convex
linear functional $\ell(\cdot)$, a linear PI inequality is a convex
optimization of the following form
\begin{align*}
& \min\limits_{P_i, Q_{1i}, Q_{2i}, R_{0i}, R_{1i}, R_{2i}}  \ell(\{P_i,
Q_{1i}, Q_{2i}, R_{0i}, R_{1i}, R_{2i}\})\notag \\
& \mrm{s. t. }\ \sum\limits_{j = 1}^K \mcl E_{ij}^*\prod \begin{bmatrix}
P_i    & Q_{1i}  \\
Q_{2i} & \{R_i\} \\
\end{bmatrix}\mcl F_{ij} +\mcl G_i \succeq 0
\end{align*}
\end{definition}
Only LMIs are required to solve an optimization problem involving LPIs.
We next describe PIE systems and their stability properties,
following~\cite{shivaknew}.
Given the PIE system $G\{\mcl{T}, \mcl{A}, \mcl{B}, \mcl{C}, \mcl{D}\}$
with PI operators $\mcl{T}, \mcl{A}, \mcl{B}, \mcl{C},$ and $\mcl{D}$ of
compatible dimensions,
and signals $w(t)\in \RL^{\mbf{n}_{\mrm{w}}}[a,b]$, we say
$\mbf{x}(t)\in\RL^{\mbf{n}_{\mrm{x}}}[a,b]$ and $z(t)\in \RL^{\mbf{n}_{\mrm{z}}}[a,b]$
satisfy the PIE with initial condition $\mbf{x}_0$ if
\begin{equation*}
\bmat{\partial_t(\mcl{T}\mbf{x})(t)\\ z(t)} = \bmat{\mcl{A} & \mcl{B} \\
\mcl{C} & \mcl{D}}\bmat{\mbf{x}(t) \\ w(t)}
\end{equation*}
for all $t\geq0$ and $\mcl{T}\mbf{x}(0) = \mcl{T}\mbf{x}_0$. Additionally, for
signals $\underline{w}(t)\in\RL^{\mbf{n}_{\mrm{z}}}[a,b]$
of compatible dimensions, we say $\underline{\mbf{x}}(t)\in\RL^{\mbf{n}_{\mrm{x}}}[a,b]$
and $\underline{z}(t)\in\RL^{\mbf{n}_{\mrm{w}}}[a,b]$
satisfy the dual PIE, $G^\top$, with initial condition $\underline{\mbf{x}}_0$ if
\begin{equation*}
\bmat{\mcl{T}^*\underline{\dot{\mbf{x}}}(t)\\ \underline{z}(t)}
=
\bmat{\mcl{A}^* & \mcl{C}^*\\
\mcl{B}^* & \mcl{D}^*}
\bmat{\underline{\mbf{x}}(t)\\ \underline{w}(t)},
\end{equation*}
for all $t\geq0$ and $\mcl{T}^*\underline{\mbf{x}}(0) =
\mcl{T}^*\underline{\mbf{x}}_0$.
For these systems, we consider input--output and exponential stability. A notational convenction
is using underlining anything related to the dual system.
Given a PIE system $G$, we say $G$ is input-output stable if for any $z, w \in L_{\mrm{e},[a,b]}$ satisfying $z = G w$ with $\mbf{x}_0=0$, we have
\[
\|p_T z\|_{L_{[a,b]}}
\leq
\gamma\|p_T w\|_{L_{[a,b]}},
\quad
\forall\,T\geq0,\quad
\forall\,w\in L_{\mrm{e},[a,b]}^{\mbf{n}_{\mrm{w}}}.
\]
Given a PIE system $G$, we say that the homogeneous part of $G$
is \emph{exponentially stable} with decay rate $\alpha>0$ if there exists
some $M>0$ such that for any $\mbf{x}_0\in\RL$,
if $\mcl{T}\mbf{x}(0) = \mcl{T}\mbf{x}_0$ and $\mcl{T}\dot{\mbf{x}}(t)
=\mcl{A}\mbf{x}(t)$, then
\[\|\mcl{T}\mbf{x}(t)\|_{\RL} \leq M\|\mbf{x}_0\|_{\RL} e^{-\alpha t} \quad\forall\, t\geq 0.\]
To describe feedback interconnections, we use the following notation.
For compatible causal $G$ and $H$, the Redheffer star
products~\cite{https://doi.org/10.1002/sapm1960391269}
denote the feedback maps
\[
\bmat{y_1\\y_2}=G\bmat{u_1\\u_2},\quad \begin{aligned}
&G\star H:u_2\mapsto y_2\\ &H\star G:u_1\mapsto y_1
\end{aligned}, \quad \begin{aligned}
&u_1=Hy_1\\ &u_2=Hy_2
\end{aligned},
\]
provided the respective interconnections are well-posed.
Here, $P\star\Delta$ closes the uncertainty loop for the generalized plant
$P$, while $K\star G$ closes the controller loop for the controlled plant
$G$.

\subsection{PIETOOLS}
PIETOOLS is a MATLAB toolbox developed by Shivakumar, Das, and
Peet~\cite{9147712} for declaring, manipulating, and solving LPIs. Its operator
syntax, PDE-to-PIE conversion, and optimization workflows are documented by
Shivakumar et al. in the user manual~\cite{pietools2025manual}. See also the
\href{https://control.asu.edu/pietools/pietools}{PIETOOLS homepage}.
